Of course, both of these relations should be understood in the sense of distributions. Thus the problem of establishing the orthogonality
and completeness relations for the eigenfunctions of the open Toda chain is equivalent to proving these two integral identities. The
problem is greatly simplified by the following remarkable result due to R.A.~Gustafson~\cite[Theorem~5.1]{MR1139492}
%
\begin{align}\label{gustafson-10}
%
\int_{\mathbb R^N}\,
\prod_{k=1}^{N+1}\prod_{j=1}^N \Gamma( \alpha_k-{\rm i}\lambda_j)\,\Gamma({\rm i}\lambda_j+\beta_k)\,
\mu_N(\lambda)\, {\rm d}^N\lambda
=\frac{\prod_{k,j=1}^{N+1}\Gamma(\alpha_k+\beta_j)}
{\Gamma\left(\sum_{j=1}^{N+1}(\alpha_j+\beta_j)\right)},
%
\end{align}
%
where $\text{Re}(\alpha_k)>0$, $\text{Re}(\beta_k)>0$ for all $k$. Note, that the integral in the l.h.s.~\eqref{gustafson-10} is exactly
the integral appearing in~\eqref{first} which, therefore, can be brought to the form
%
\begin{align}\label{first-reduced}
\lim_{\varepsilon,\varepsilon' \rightarrow 0^+} (2\pi) \delta(\Lambda-\Lambda')\,
\frac{\prod_{k,j=1}^{N+1}\Gamma({\rm i}\lambda'_j-{\rm i}\lambda_k+\varepsilon+\varepsilon')}
{\Gamma\big({\rm i}\Lambda'-{\rm i}\Lambda+N(\varepsilon+\varepsilon')\big)} = {\mu_{N+1}^{-1}(\lambda)}\,
\delta^{N+1}\left(\lambda,\lambda'\right).
\end{align}
The proof of this identity is already rather straightforward. Some details can be found in~Appen\-dix~\ref{Appendix identite fct delat
evoluee}.

It takes a little more work to prove the identity~\eqref{second}. Having put $\alpha_{N+1}=L$ and $\beta_{N+1}=t L$, $t>0$,
in~\eqref{gustafson-10} and sending $L\to +\infty$ one arrives at the reduced version of the integral~\eqref{gustafson-10} which takes the
form~\cite{DMspinchainandSL2RGustafson}
\begin{gather}\label{gustafson-1a}
\int_{\mathbb R^N} t^{{\rm i}\Lambda}
\prod_{k,\,j=1}^{N} \Gamma(\alpha_k-{\rm i}\lambda_j)\,\Gamma({\rm i}\lambda_j+\beta_k)\,
\mu_N(\lambda)\, {\rm d}^N\lambda
=\frac{t^{A}}{(1+t)^{A+B}}{\prod_{k,\,j=1}^{N}\Gamma(\alpha_k+\beta_j)},
\end{gather}
where $A(B)=\sum_{k=1}^N \alpha_k(\beta_k)$. At the same time representing the l.h.s.\ of equation~\eqref{second} in the following form (for
more details see Lemma~\ref{Lemme isometrie operateur SAB}):
\begin{gather*}
\lim_{L\rightarrow\infty} \lim_{\varepsilon,\,\varepsilon'\rightarrow 0^+} \frac1{\Gamma^{{2N+2}}(L)}\int_{\mathbb R^{N+1}}
\prod_{k,j=1}^{N+1} \Gamma({\rm i} \gamma_k-{\rm i}\lambda_j+\varepsilon)\,
\Gamma({\rm i}\lambda_j-{\rm i}\gamma'_k+\varepsilon')
\\ \hphantom{\lim_{L\rightarrow\infty} \lim_{\varepsilon,\,\varepsilon'\rightarrow 0^+} \frac1{\Gamma^{{2N+2}}(L)}\int_{\mathbb R^{N+1}}
\prod_{k,j=1}^{N+1}}
{}\times{\rm e}^{{\rm i}\Lambda(x'_N-x_N)}
\mu_{N+1}(\lambda) \,{\rm d}^{N+1}\lambda,
\end{gather*}
where ${\rm i}\gamma_{N+1}=-{\rm i}\gamma'_{N+1} =L$, one can evaluate the integral using equation~\eqref{gustafson-1a}. Then, after some algebra, one can
show that \eqref{second} is equivalent to the following identity
\begin{gather}
\lim_{L\rightarrow\infty} \lim_{\epsilon\rightarrow 0^+}
L^{{\rm i}(\Gamma'-\Gamma)}\prod_{k,j=1}^{N}\Gamma\big({\rm i}(\gamma'_k-\gamma_j)+\varepsilon\big)
\sqrt{\frac{L}{4\pi}}{\cosh^{-2L}\bigg(\frac{x_N-x'_N}2\bigg)}\nonumber
 \\ \hphantom{\lim_{L\rightarrow\infty} \lim_{\epsilon\rightarrow 0^+}}
{}= {\mu^{-1}_{N}(\gamma)}
\,\delta(x_N-x'_N)\,\delta^{N}(\gamma,\gamma').\label{second-reduced}
\end{gather}
We recall here that equations~\eqref{first-reduced} and \eqref{second-reduced} have to be understood in the sense of distributions and relegate
further details to Appendix~\ref{Appendix identite fct delat evoluee}.
